% Version 3.10: persistent regional request ownership and finite hydration.
% Included inside Requests, after request definitions and public import.
\subsubsection{Persistent request owners and roster hydration}
\label{sec:ownership}
\label{sec:request-ownership}
A node that receives a still-pending request keeps a node-level request
service. Closing its production program does not close that service.
Before a new roster prepares its input selection, it obtains finite cuts
from historical owners in the same physical region and disseminates the
acquired bytes to its own members. The procedure below defines the
directory, cuts, transfers and completion certificate. It neither
publishes bare requests nor changes the original complete result $R_0$.

\paragraph{Directory and communication scope.}
Write $V_i(k)$ for region $i$'s membership in $\tagg{AUTH}_{k,0}$ and set
\[
 D_i(r)=\bigcup_{0\le k\le r}V_i(k),\qquad
 d_i(r)=\max_{0\le k\le r}\Delta_i(k).
\]
Deduplicate the union by authenticated persistent identity, retaining
the historical membership evidence. The maximum is positive and finite.
Here $r$ is an actually admitted target slot. Genesis starts at zero;
regular succession and a Force after positive accepted output advance
the responsible slot by one; a zero-output retry keeps that slot.
Induction on these transitions shows that admission of $r$ has a finite
predecessor lineage containing the immutable AUTH slots $0,\ldots,r$.
A Force may replace speculative successors, but does not erase the
authority records that established these slots. A later authorized
checkpoint retains their identities and authentication dependencies.
Obtain missing records by their known keys $(k,0)$ from this lineage.
This is finite known-key retrieval, not enumeration of all L1
authorities, a search for absent indices, or a new directory oracle.
Acquisition of this fixed evidence is preparation with no new deadline.

The within-region guarantee in \cref{ass:network} covers correct old
and new AUTH members of this same physical region while their protocol
duties continue. The exchanges below therefore use the existing
allowances in the finite prefix, bounded by $d_i(r)$. This is no
cross-region bound and takes no maximum over future configurations.
An identity whose pending and outgoing duties have ended need not
remain online merely because it occurs in $D_i(r)$. Its silence is
handled by a finite timeout. A correct owner whose pending bytes are
needed remains responsible for their service as specified below.

Fix a positive chunk payload $b$ within the original admitted
complete-message limits. Each envelope supplies its authenticated
sender, type, chain, target context, exact object descriptor and index.
The original size, complete-validation and aggregate-load scope applies
to these actual envelopes and dispatches. A finite pool is sent in
finitely many such envelopes; its arbitrary size is never charged to
one regional delay. The timed source-authentication handler below is
the existing supplied-dependency L1/AUTH check. It does not perform
missing-history acquisition, application replay or native restoration.
Those operations, and any whole-object computation outside the original
complete handler, finish as actual preparation. No new deadline or
larger service envelope is assigned to them.

\paragraph{Persistent state and finite cuts.}
Every correct member keeps the following state outside
$(\beta,r,g,t)$ and preserves it across worker exit and restart.
\begin{itemize}[itemsep=2pt]
\item $\tagg{Pending}$ retains complete authenticated request bytes and
their logical identities. Only verified inclusion in $G$ permits
removal. An acknowledgement or another roster's copy does not.
\item $\tagg{ReceiptJournal}$ is an append-only sequence of first
complete authenticated request receipts, with a maintained finite
length, prefix commitment and indexed records. A duplicate changes
neither its first-receipt position nor its identity.
\item $\tagg{RawJournal}$ records complete immutable raw-object
handles. A handle names the exact stored fragments, finite length,
digest and authentication material. Processing the last required
fragment appends the handle using incremental bookkeeping; it does not
run canonicalization, application admission or a whole-pool scan.
\item $\tagg{Cuts}[i,r,v]$ stores the immutable finite journal cut
promised to target member $v$, its authenticated roster, and unfinished
indexed replies. $\tagg{Targets}$ stores rosters registered by valid
queries. Context records, transfer cursors, shares, collectors and
unfinished sends are also durable.
\end{itemize}
The journal starts finite. A complete first-receipt handler appends at
most one finite entry; every other transition preserves a finite
sequence or takes an actual finite prefix. Induction on the local
transitions therefore gives a finite length $J_p$ at every reached
query handler. There is no limit transition that installs an infinite
journal. Taking a cut copies the maintained length/root, not a scan.
This local invariant does not assert a finite number of events at all
nodes before a given clock time. The regional complete-message bound
also prevents a required correct query handler from being postponed
behind an infinite unfinished local prefix.

On a first complete authenticated request receipt, persist its bytes
and journal entry before acknowledging receipt. Dispatch its complete
bytes to the node's installed rosters and registered target rosters.
Receiving nodes retain the bytes by the same rule. Deduplicate a first
multicast by request identity and target AUTH; keep unfinished sends.
New complete raw records remain in $\tagg{RawJournal}$ until their
actual canonical processing in \cref{sec:input-selection}; duplicates
also retain a complete raw or canonical handle. Later invalidity at one
speculative parent does not delete a request. Requests already in $G$
may be marked unnecessary for a new cut, but bytes pinned by evidence,
selection or outgoing work remain available.

The original correct owner retains a pending request even after every
known target acknowledges it. Future rosters consequently query the
same authenticated owner; they do not chase a moving private custodian.
An identity first appearing in slot $k$ is included in later directories
$D_i(r)$ with $r\ge k$. No earlier roster is claimed to discover private
requests at identities outside its authenticated prefix.

\paragraph{One preparation per member and slot.}
Choose protocol parameters $\tau_p>0$ and common clock origin $\tau_0$.
Define
\[
 \operatorname{release}(r)=\tau_0+r\tau_p.
\]
For each $(i,r,v)$ there is one immutable preparation and first-query
event. A correct $v$ starts that query only at a clock time
$s_v\ge\operatorname{release}(r)$. Earlier context preparation is
allowed. Beta, recovery and changes of production worker do not change
the key or restart the job. This is a local protocol wait, not a
minimum network delay or a global event-rate premise. Its deadline is
finite for every fixed $r$. It does not pause or reset public windows,
Force accounting, source timestamps, Ready entry times or post-entry
timers.

Let $N=V_i(r)$ and $d=d_i(r)$ for the rest of this procedure. A
preparation is authenticated as the named member $v\in N$; another
sender cannot create $v$'s cut by writing $v$ in a packet.

\paragraph{Complete context before the first query.}
For each $p\in D_i(r)$, the correct preparer $v$ executes these actions.
\begin{enumerate}[label=(\arabic*),itemsep=3pt]
\item Construct and validate a finite authentication capsule $A_{v,p}$
from the existing trusted chain checkpoint. It supplies the required
finalized L1 headers and proofs, complete applicable control evidence,
the immutable AUTH prefix through $r$, historical identities, and the
exact target context. Every referenced validation dependency is actual
supplied content. Retrieval and local preparation finish before the
timed schedule; an unresolved archive pointer is not a capsule.
\item Freeze its exact finite length $L$ and digest $H$, with
$n=\lceil L/b\rceil$ fixed once. Actually send its metadata and all
$n$ indexed chunks to $p$. Raw authenticated storage is permitted
before $p$ knows the proposed target AUTH; sender and context keys
prevent substitution. Let $a$ be the time of the last actual envelope
send to any directory member. Merely queuing a send does not define
$a$. The finite dispatch waits for no recipient acknowledgement.
\item At or after $a+2d$, actually send every $p$ a complete
$\tagg{CONTEXT\_VALIDATE}(i,r,v,p,H)$ trigger. Let $a'$ be the last
actual trigger-send time. By $a+d$, each required correct owner has
the complete cached capsule, strictly before these triggers. Processing
the trigger runs only the original complete source-authentication
predicate on those supplied dependencies. It stores
$\tagg{ContextReady}[i,r,v,p,H]$ durably. Invalid or genuinely missing
dependencies do not grant that record or silently defer the same
timed success predicate to a background job.
\item Begin the first $\tagg{CUT\_QUERY}$ multicast at actual time
\[
 s_v\ge\max\{a'+2d,\operatorname{release}(r)\}.
\]
Send to every directory identity without waiting for replies. Record
each actual send time $s_{v,p}\ge s_v$ separately. Every required
correct $p$ has processed its validation trigger by $a'+d<s_v$.
Thus each query starts with complete authenticated receiver context;
the sender needs neither the correct subset nor an all-member ACK.
\end{enumerate}
The strict guards handle equal clock timestamps without ordering
different nodes' handlers. They are conservative local waits using the
existing delay. Authentication of a supplied complete capsule is not a
bound on arbitrary from-genesis catch-up or native-state conversion.
Such work cannot be inserted into the timed trigger. Cached context,
actual send times and unfinished duties survive every beta change.

\paragraph{Immutable descriptors and record slots.}
A query binds $(i,r,v)$ and $p$'s capsule digest. The receiver rejects
an early query before $\operatorname{release}(r)$ without registering
the target or reserving a cut; it does not defer that packet for later
authorization. For the first valid query, $p$ checks the authenticated
sender and its $\tagg{ContextReady}$ record, then atomically registers
$N$ and freezes its current $m=J_p$ and prefix root. It immediately
returns a signed descriptor $(i,r,v,p,m,\mathit{root})$. Repeated
queries reuse that descriptor. New receipts append only to the live
journal and use the continuing push path.

At $s_{v,p}+2d$, after messages due at that time, $v$ freezes at most
one complete valid descriptor from $p$. Use the first accepted one if
a Byzantine owner equivocates. Missing or malformed replies give no
descriptor. Later roots or larger counts cannot replace it. Each
correct needed owner processes the query within $d$ and returns its
complete descriptor within the second $d$. The finite directory yields
finitely many deadlines, without an identifiable ``all correct replies
received'' event.

Process the frozen descriptors and their finite record indices in
deterministic local order. Write
$C=(\mathit{chain},i,r,v,p,m,\mathit{root})$ for a fixed descriptor.
For each $0\le j<m$ perform this procedure exactly once.
\begin{enumerate}[label=(\arabic*),itemsep=3pt]
\item Actually send $\tagg{RECORD\_META\_QUERY}(C,j)$ and allow one
$2d$ query/reply slot. The reply binds $C,j$ and supplies the exact
finite record length $L\ge0$ and digest $H$, or a signed tombstone for
an entry already included in $G$. Freeze the first matching complete
valid metadata by the deadline. Missing or malformed metadata, or a
tombstone, ends this record without adding a request. Metadata changes
never restart it. A tombstone proves no new global acceptance.
\item Fix $n=\lceil L/b\rceil$ once. If $L=0$, take no fragment
slots and proceed to final checking. Otherwise, for each
$h=0,\ldots,n-1$, actually send
$\tagg{FRAGMENT\_QUERY}(C,j,L,H,h)$ and allow exactly one $2d$ slot.
Its signed reply binds these exact fields. Require $b$ payload bytes,
except that the final payload has $L-b(n-1)$ bytes. Retain at most one
matching valid payload per index. Process due replies before expiry.
\item Ignore unauthenticated, wrong-sender, wrong-context and
out-of-slot packets. They cannot fail a correct owner's record. A
malformed or conflicting reply signed by the expected owner for this
active context may fail that record; absence of a matching reply at
expiry also fails it. Stop or consume the remaining fixed slots on
failure, but never add slots, enlarge $L$, replace $H$ or restart.
\item For successful reconstruction, check the exact $L$ bytes and
digest $H$, canonical encoding and request authentication. Perform
this actual terminating local work before adding a request to the
acquired pool. No missing external dependency is awaited and no new
$d$ deadline is assigned to this whole-object work. Application
admission after a future exact parent is left to input selection.
\end{enumerate}
Thus each record consumes at most one metadata slot and $n$ fragment
slots plus its actual finite local checks. Byzantine finite counts and
lengths may be large, but cannot grow a job after it starts. This is
not a practical storage or denial-of-service bound. Every correct
needed owner supplies the immutable bytes within each applicable slot.
If the request enters $G$ during transfer, a correct tombstone is
harmless. Otherwise it reconstructs and passes the local checks.

All slot progression, indexed replies and sends are protocol actions
with durable cursors. On completing or timing out one step, continue
the next fixed step without a Byzantine retry, new beta or unsolicited
permission. Continuing local execution and the original admitted load
scope are used here; no external fair-ingress service is added. At the
end freeze $\mathit{Pool}_v$. Live arrivals do not enlarge it. It
contains every request still absent from $G$ that was held at a correct
needed owner's query cut, including every such request held before
$s_v$. Finish any remaining local checks before dissemination.

\paragraph{Actual all-member dissemination and HYDRATED.}
For each frozen pool record, $v$ actually sends the metadata and all
fixed chunks to every member of $N$. These are finite multicasts of
already held content, without a query phase. Recipients stage out-of-order
fragments and deduplicate by sender, target, record and chunk index.
Incremental storage exposes the complete raw handle in
$\tagg{RawJournal}$ when all fixed fragments are present. This does
not require a whole-record join, canonicalization or application scan
in the last-chunk handler. Those later computations read the retained
handle during actual preparation.

The node-level receiver runs without a production engine and while
its own acquisition is unfinished. It also runs after a worker exits.
Let $b_v$ be the time of the last actual envelope send to any member,
including every metadata and chunk envelope. For an empty pool, use
the time acquisition and local preparation finish. Only at or after
$b_v+2d$ may $v$ create its signed share on the common context
\[
 \mathit{HC}_i(r)=
 (\mathit{chain},\tagg{HYDRATED},i,r,
   \Hash(\tagg{AUTH}_{r,0}),\mathit{AP}_r),
\]
where $\mathit{AP}_r$ commits to the canonical ordered immutable AUTH
records $0,\ldots,r$. Alternate proof encodings do not change it.
The context contains no beta or pool root. Different correct
preparers may acquire different supersets. All required raw handles
are retained at each correct target by $b_v+d<b_v+2d$, strictly before
the share. This is byte availability, not a claim of application
eligibility after an unknown parent.

Share and certificate dispatch are mandatory actions:
\begin{enumerate}[label=(\arabic*),itemsep=3pt]
\item At share creation, persist the full signed share, insert it into
the local collector and actually multicast it to every target member.
Unfinished per-recipient sends remain duties across beta replacement,
storage handoff and worker exit.
\item Each target manager checks the share's signature, membership,
exact $\mathit{HC}_i(r)$ and fixed target voting snapshot. Retain at
most one valid matching share per distinct signer. Duplicates add no
weight; another context has a separate collector.
\item As soon as collected weight exceeds $F_i(r)$, materialize and
persist the full $\tagg{HYDRATED}$ certificate. On first possession
of a complete valid certificate, actually multicast it to all target
members. Each receiver checks the complete signature set and weight,
persists it and performs the same once-per-context relay. A root or
unresolved pointer is insufficient. If fragmented before use, its
exact finite length and chunk count are fixed by the same procedure.
\item Make the complete checked certificate available to the
preparation manager. Continue all promised storage, indexed replies,
share sends and certificate relays. No bare share or certificate gains
a $T_{\rm pub}$ or $T_E$ guarantee.
\end{enumerate}
A member may adopt a complete valid certificate and cancel redundant
private acquisition, preserving every promised service and outgoing
duty. It need not identify the correct signer or fetch that signer's
pool after adoption. Without such a certificate, every correct target
continues its finite preparation and sends its share.

\begin{proposition}[Finite fixed-slot hydration]
\label{prop:ownership-formation}
For a fixed admitted $(i,r)$, after its finite required AUTH/context
evidence is available, correct target members complete hydration in
finite time under the original communication and processing scope.
This time may depend on the actual finite cuts and preparation work;
it is not a uniform archive, native-restoration or future-roster bound.
\end{proposition}
\begin{proof}
The directory and each capsule are finite. Their actual finite
preparation and dispatch complete. The context barriers and
$\operatorname{release}(r)$ have finite deadlines. Every descriptor
deadline expires, and each retained descriptor fixes finitely many
finite record jobs. Malformed or silent owners cannot extend them.
Each local check and fixed multicast completes, so every correct
preparer either adopts an existing certificate or reaches and sends
its share. If none exists earlier, take the finite maximum of the
correct members' share-send and delivery times. Every correct collector
then has correct weight
$W_i(r)-B_i(r)\ge W_i(r)-F_i(r)>F_i(r)$ and forms a certificate.
Adoption only makes this earlier. No Byzantine contribution is needed.
\end{proof}

\begin{proposition}[Hydration precedes every certificate]
\label{prop:ownership-hydration}
Fix any valid $\tagg{HYDRATED}$ certificate for $(i,r)$ and any one
of its correct signers $v$. Every request pending at a correct
directory owner before $s_v$ and still absent from $G$ has a complete
retained raw or canonical handle at every correct target member
strictly before this certificate can exist.
\end{proposition}
\begin{proof}
The correct owner is a continuing source while that request is pending.
The completed context barrier makes its first query valid. Its exact
cut contains the request; each fixed metadata/fragment reply arrives
within its slot. Hence $v$ checks and includes the bytes in its frozen
pool. Its actual multicast puts the complete handle at every correct
target by $b_v+d$. This strictly precedes its share and therefore every
certificate containing that share. Every valid certificate has a
correct signer because its weight exceeds the faulty-weight bound.
No intersection with a later offer or REPORT signer set is needed.
\end{proof}

\paragraph{Input selection and Ready.}
The certificate is a precondition for a correct member's finite
selection cut in \cref{sec:input-selection}. That cut includes raw
handles whose canonicalization is still pending. The resulting
$\tagg{OPENING\_OFFER}$ and opening certificate carry the selection
claim. Ready transmits that opening certificate, as defined there;
it does not additionally transmit the HYDRATED signature set or the
historical directory capsule. A correct offer signer already completed
that causal preparation. Adoption requires neither a private pool
fetch nor completion of the adopter's own redundant hydration job.
The original executable-state and complete Ready-message interfaces
remain in force.

For fixed $r$ and anchor, include this finite hydration completion in
the preparation maximum. Its historical $d_i(r)$ is spent before
Ready. The target snapshot's $\Delta_i(r)$ still governs the original
timed regional messages. No whole-pool transfer is inserted into a
post-Enter response constant. A changed beta alone reuses hydration;
an exact-anchor selection change is handled by the separate selection
key. Later live pushes never reopen a cut, amend a share or reset a
timer.

\begin{proposition}[Finite early target slots]
\label{prop:ownership-finite-exceptions}
Suppose a correct owner first receives a fixed request $x$ at finite
clock time $t_x$. Only finitely many target slots can have a correct
first query at or before $t_x$.
\end{proposition}
\begin{proof}
For any such query,
$\tau_0+r\tau_p\le s_v\le t_x$. Its slot lies in
\[
 E_x=\{r\in\mathbb N_0:\tau_0+r\tau_p\le t_x\}.
\]
This set is finite since $\tau_p>0$; it is empty if $t_x<\tau_0$.
Equality is included, so no equal-time receipt/query order is needed.
Each authenticated target roster is finite and each $(i,r,v)$ has one
first query. Summing its finite size over authenticated targets in
$E_x$ bounds the number of early correct first queries; an index with
no authority or target contributes none. Byzantine attempts cannot
create a correct member's cut.
Activity at unrelated nodes does not enlarge this set.
\end{proof}

\begin{proposition}[Coverage across future rosters]
\label{prop:ownership-coverage}
Let a correct owner $p\in V_i(k)$ receive $x$ at $t_x$ and retain its
service while $x$ is absent from $G$. For every target slot
$r\notin E_x\cup\{0,\ldots,k-1\}$, every valid HYDRATED certificate
for $(i,r)$ has the following property: complete $x$ bytes and a raw
or canonical handle are retained at every correct target member
strictly before that certificate exists, unless $x$ has already
entered $G$. Only finitely many accepted DATA contributions of region
$i$ can belong to the excluded target slots.
\end{proposition}
\begin{proof}
For a nonexcluded target, $p\in D_i(r)$ and every correct signer's
first query is strictly after $t_x$. Apply
\cref{prop:ownership-hydration} to a correct signer in any proposed
certificate. This proves the universal certificate claim, including
a certificate assembled later from old shares. A correct old share
whose query preceded $t_x$ could only belong to $E_x$.

The excluded slots form a finite set. Each has a fixed finite DATA
quota. Zero-output beta retries contribute no accepted DATA and reuse
the slot. A positive-output Force or regular completion advances the
responsible slot; discarded speculative successors cannot contribute
accepted DATA after closure. Thus an excluded numeric slot cannot
multiply its accepted quota through beta retries. A finite set of such
slots contributes only finitely many accepted DATA. Consequently,
infinitely many accepted contributions from $i$ eventually use covered
slots, without roster overlap, a finite union of future identities, a
uniform future delay or a global event-rate assumption.
\end{proof}

\paragraph{Worker retirement and physical exit.}
The manager-level owner keeps pending bytes, exact cuts, raw and
canonical handles, context records, share collectors, fragments and
unfinished sends in every program phase. Closed,
$\tagg{RequestClose}$ and Terminated do not disable these handlers.
A worker may be released after the retirement guards transfer its
storage, receive route and duties to that persistent service. It may
not release the last required bytes or replace an actual owner by an
unresolved retrieval task. Existing shares and certificates may still
be served and relayed; these actions never create a regional vote,
restart an engine or reopen a closed fresh-signature gate.

A physical machine may cease this request service only when it holds
no still-pending request whose original ownership it must preserve and
has no remaining pinned or outgoing duty, or when the same logical
owner's complete state, authenticated receive route and continuing
service in the same physical region already survive the local
worker/machine handoff. An arbitrary
new custodian's acknowledgement is not that handoff. In particular,
the last correct original owner of a still-pending request cannot
leave merely because its AUTH is old, the successor is Ready, or a
HYDRATED certificate exists. The protocol promises worker retirement,
not an independent deadline for physical departure or bounded storage.
This is the responsibility of correct participants under the existing
state-retention and communication model; no new archival service or
external party is introduced.
